\documentclass[12pt]{article}
 \usepackage[german]{babel}
 \usepackage[utf8]{inputenc}
 \usepackage{graphicx}
 \author{Denis Herdt, Almin Causevic}
 \title{ROS auf dem Rapberry Pi}
 \setlength{\parindent}{0pt}                   % Einrueckung 1. Zeile eines Absatzes
 \setlength{\parskip}{5pt plus 2pt minus 1pt}  % Abstand zwischen Absaetzen
 \frenchspacing
 \sloppy
 
 
\begin{document}
%4-20 Seiten
\begin{figure}[h]

\includegraphics[width=4cm]{hs-logo.jpg}
\end{figure}
Fakultät für Elektrotechnik und Informatik

\vspace{3cm}

\begin{center}

{\bf \huge ROS auf dem Raspberry Pi}
\vspace{4cm}

14 November, 2013
\vspace{1cm}

Systemadministration Projekt in Angewandter Informatik \\
von Denis Herdt und Almin Causevic

\end{center}

\pagebreak

\tableofcontents

\pagebreak

\section{Einleitung}
1-2 Seiten
\subsection{Motivation}

Wir haben uns für dieses Thema entschieden, weil wir in Kombination mit Hard- und Software arbeiten wollten. Wir haben mitbekommen, wie umständlich z.B. der Roboter ''Volksbot'' im Robotiklabor der Hochschule mit einem zusätzlichen, auf dem Roboter platzierten Laptop gesteuert werden muss, der unnötig Platz und Gewicht beansprucht. Deshalb haben wir uns entschlossen, dieses Problem mithilfe des kleinen und leichten Raspberry Pi zu lösen.
Die Idee stieß auch auf großes Interesse bei den Labormitarbeitern und bringt viel praktischen Nutzen.

Unser Ziel möchten wir mithilfe des Software-Frameworks ROS (Robot Operating System) realisieren.
Für ROS haben wir uns entschieden, weil es zur Zeit die aktuellste und inovativste Methode ist, ein flexibles Netzwerk von Komponenten für z.B. die Robotersteuerung zu realisieren. Es wird weltweit eingesetzt und könnte auch in Zukunft sehr nützlich werden.

Durch die Kombination mit der Mächtigkeit und Flexibilität von ROS und der geringen Größe und dem niedrigen Energieverbrauch des Pi eröffnen sich außerdem neue und effizientere Einsatzmöglichkeiten.

Sehr gut finden wir auch die Tatsache, dass wir mit Linux arbeiten können und dieses
%Konfiguration von Linux???
Projekt viel Konfiguration des Betriebssystems voraussetzt, da wir auch gerne mehr
%du meinst Praxiserfahrung??
Praxisanwendung mit Linux haben möchten.

%Roboter Lab Anwendung von ROS und Möglichkeit, Roboter zu nutzen

%Interesse der Lab-Crew an Umsetzung ROS mit Raspberry Pi

%Kombination Hard-/Software

%Einarbeitung in Kommunikations Framework ROS

%Mit Hardware in Berührung zu kommen

%Konfiguration mit Linux

%Nodes nun mit Hardware (Pi) realisierbar

\subsection{Zielstzung}


Unser Hauptziel besteht darin, ROS auf dem Raspberry Pi aufzusetzen und zu nutzen. Dieses Ziel wird wie folgt unterteilt:\\

Zunächst möchten wir ein passendes, auf Linux basiertes Betriebssystem auswählen und auf dem Raspberry Pi aufsetzen. In diesem Fall haben wir uns für Raspbian entschieden, da es außerdem am kopatibelsten zu ROS ist.

Als nächstes wählen wir einen geeigneten Roboter aus. Die Kriterien dafür erarbeiten wir uns aus einer vorangegangenen ROS-Recherche.

Im Anschluss setzen wir ROS auf dem Raspberry Pi und mehreren Linux-Rechnern auf und stellen ein WLAN-Netzwerk zwischen den verschiedenen ROS-Komponenten her.
Dieses Netzwerk wird mithilfe von Testausgaben überprüft.

Damit ist die Grundlage zum Anwenden von ROS auf dem Pi geschaffen, an die man weiter Projekte knüpfen kann.\\

{\bf optional}\\

Bleibt am Ende noch genug Zeit, implementieren wir die Steuerung des ''Volksbots'' auf dem Raspberry. Zur Usabilitysteigerung integrieren wir ein Smartphone ins ROS-Netzwerk. Die Bewegungsrichtung und -geschwindigkeit des Roboters soll durch das Schwenken des Smartphones bestimmt werden.

Haben wir damit Erfolg und wiederum genügend Zeit zur Verfügung, sprechen wir die integrierte Kamera des Roboters an und streamen das Bild mithilfe des Netzwerkes auf einen externen Bildschirm.

Mit diesen Optionen wollen wir zeigen, wie Flexibel und Umfangreich ROS in der Anwendung sein kann.

%Roboter durch einen Raspberry Pi mithilfe von ROS steuern
%Bild einer Webcam durch ROS auf einen Bildschirm streamen

\subsection{Eigene Leistung}

Das ''neue'' an diesem Projekt bezieht sich
%Unsere Leistung bezieht sich
auf die Ablösung eines Laptops durch den kleineren und leichteren Raspberry Pi als ROS-Komponente am ausgewählten Roboter, sowie einem Leistungs- bzw. Eignungstest des Raspberry 
im Aufgabenbereich Robotersteuerung.
% als Steuerungseinheit des "Volksbot".


%anstatt Linux PC wird Pi benutzt

\subsection{Aufbau der Arbeit}

\begin{itemize}

\item Recherche über Linux Systeme auf dem Raspberry Pi
\item Passendes Linux System auf Raspberry Pi gesetzt			
\item Recherche über ROS
\item Geeigneten Roboter gewählt (Recherche Verfügbarkeit, Komponenten etc.)
\item Komponenten geschafft (Wlan Stick, Akkupack, etc.)
\item ROS auf Raspberry Pi aufgesetzt
\item Netzwerk zwischen mehreren ROS-Komponenten hergestellt
\item Netzwerk mithilfe von Testausgaben überprüft \\

\end{itemize}

{\bf optional}

\begin{itemize}

\item Robotersteuerung wird zur Zeit implementiert


%Roboter durch einen Raspberry Pi mithilfe von ROS steuern
%Bild einer Webcam durch ROS auf einen Bildschirm streamen

\end{itemize}

\section{Grundlagen}
%6-8 Seiten

\subsection{Raspberry Pi Hardwarekomponenten}

Der Raspberry Pi ist ein kreditkartengroßer Einplatinencomputer, der von der Raspberry Pi Foundation entwickelt wurde.
Wir benutzen für dieses Projekt das leistungsstärkere Modell B.

Technische Details:

\begin{itemize}
\item Preis: ca. 35 Euro
\item Prozessor: ARM1176JZF-S (700 MHz)
\item Broadcom VideoCore IV
\item SDRAM: 512 MB
\item Bis zu 16 GPIO-Pins
\item USB-Anschlüsse: 2
\item FBAS, HDMI
\item 3,5-mm-Klinkenstecker (analog), HDMI (digital)
\item Kartenleser für SD (SDHC und SDXC)/MMC/SDIO
\item 10/100-MBit-Ethernet-Controller 
\item 5 V, 700 mA (3,5 Watt)
\item 5-V-Micro-USB-Anschluss (Micro-B), alternativ 4 x AA-Batterien
\end{itemize}

Für den Anschluss an einen Bildschirm benutzen den HDMI Ausgang.
Diesen benutzen wir für Kontrollausgaben und Tests.

Als Stromquelle steht ein Akkupack ......(Firma,Modell) zum Einsatz. 
Hierbei ist wichtig, dass der Raspberry mind. 700 mA, besser 1 A zur Stromversorgung bekommt. 
Bei niedrigerer Amperzahl arbeitet der PC oft nicht zuverlässig.

Ein schneller und großer RAM Speicher ist für unsere Zwecke wichtig, da viele Signale, teils sogar synchron, verarbEin 17-Zoll Monitor wird für die Bildschirmausgabe des Raspberry über HDMI zu DVI benutzteitet werden müssen.

Für die Datenübertragung über das Netz benutzen wir einen Standard 300Mbit/s Wlan-Stick.

Der Raspberry Pi arbeitet mit einer ARM-Prozessorarchitektur.
Diese Architektur wird gerne für embedded Systems, wie PDAs oder Router, eingesetzt.

Sie ist auch auf jedem Smartphone zu finden, da sie den Vorteil einer sehr geringen Leistungsaufnahme bietet. 

Erwähnenswert ist die Architektur deshalb, weil wir im Laufe des Projektes Schwierigkeiten hatten, auf die wir später eingehen werden
Wir benutzen das auf Linux basierende Betriebssystem Raspbian.
Dabei handelt sich um ein für Raspberry Pi optimiertes open-source Debian-System.
Es enthält viele für die ARM-Architektur vorkompilierte Pakete (über 35.000), dazu auch Features wie etwa eine GUI.

Das System braucht 3GB Speicherbedarf unserer 16GB großen SD-Karte.
Der zusätzliche Speicherplatz wäre nötig, falls man vorhat, Log-Dateien und Ähnliches direkt auf dem Raspberry Pi zu speichern. 
In unserem Fall übernehmen das jedoch die leistungsstärkeren PC's über das ROS Netzwerk.

GPIO (General Purpose Input/Output) ist ein weiteres interessantes Feature des Raspberry Pi.
Es gibt uns die Möglichkeit, jegliche Hardware-Funktionalität anzusteuern.
Beispielsweise können LED-Leuchten oder der Start-Knopf der Kaffeemaschine damit über ROS gesteuert werden.
Wir konnten uns leider jedoch zeitlich bedingt nicht mehr mit dieser Thematik beschäftigen, es bietet aber Anreiz für noch mehr Ansätze und Umsetzungen mithilfe des Raspberry Pi und ROS.


\subsection{Smartphone}
%App-Steffen

\subsection{Volksbot}

\subsection{ROS}
%catkin, Netwerk



\section{Problem}

Hauptproblem??

\subsection{Verwandte Arbeiten}

Steffen, Marc
ROS-Tut

\section{Anforderungen}

\begin{itemize}
\item Datenpakete über das ROS-Netzwerk verschicken
\item Beliebige Datentypen verarbeiten
\item Beliebige Linux-Systeme einbinden
\item Überprüfung von Log-Dateien, Kontrollausgaben etc..
\item Robuste und zuverlässige Kommunikation
\item Modulares und flexibles Netzwerk
\item Repositories sollen leicht in Projekte eingebunden werden
\end{itemize}

optional:

\begin{itemize}
\item Roboter soll bei Steuerung in richtige Richtungen fahren
\item Roboter soll sofort auf Richtungsanweisungen reagieren
\item Sensibilität der Steuerung sollte einstellbar sein
\item Steuerung über WLAN
\item Raspberry Pi auf Roboter ohne Kabelzgewirr %haben wir mit WLAN nicht
\item Node des Raspberry beim Hochfahren automatisch starte
\item Leichte Bedienbarkeit des Roboters
\vspace{0,6cm}
\item Streaming-Bild ruckelfrei verarbeiten
\item unkonvertiert Echtzeitübertragung
\item Ausgabe am externem Bildschirm
\end{itemize}


\section{Lösungsvorschläge}


Es wird ein ROS-Master aufgesetzt. Dieser sorgt für ein bestehendes ROS-Netzwerk und ist für die Komponentensynchronisation und Paketverwaltung zuständig.
Über Topics und Nodes werden Datenpakete verschickt und empfangen.
%TODO wiki nachschauen obs stimmt

Die Verarbeitung von verschiedenen Datentypen ist durch msg und srv realisierbar
%TODO msg/srv näher erklären

Die Unterstützung beliebiger Linuxsysteme wird von ROS grundsätzlich unterstützt.
%TODO umschreiben

Kontrollausgaben und ähnliches werden durch Horchen an Topics mithilfe von echo und durch Ausgabenbereiche im Quellcode durch das Netzwerk realisiert.

 
%Horchen an Topics mit echo + Ausgabe im Code mithilfe ROS-Stream Fkt.

Die Netzwerkstabilität und -zuverläsigkeit könnte durch Begrenzung der zeitgleichen Datenübertragung umgesetzt werden. Zusätzlich kann man eine Überprüfung der verschickten und erwarteten Daten einbinden.


%Verschickte und erwartete Datentypen müssen stimmen
%Richtige Konfiguration des Masters und Clients
Ein Modulares und flexibles Netzwerk, sowie das einfache Einbinden von Repositories ist gernerell durch ROS in kombination mit Catkin gegeben.


%durch ROS und catkin gegeben 
%durch catkin gegeben

opt:
abhängig vom richtigen Code der Motorsteuerung

abhängig von Raspberry Pi Hardware und effizientem Code und
guter Hardware

abhängig vom Code

Hardware nötig Wlan-Stick, Akku-Pack 

sollte in Linux beim booten konfiguriert werden

gegeben durch gute Smartphone App

opt:
abhängig von Raspberry Pi Hardware und Übertragungs-Codierung

abhängig vom Codex Hardware Komponenten

nötige Hardware und Netzwerk legen

\section{Bewertung der Lösungen}

alles gut

\section{Implementation}

ROS Netzwerk (topics, nodes, catkin...)
Linux SD Karte Pi
Datenübertragung an Volksbot
Konfiguration Motorsteuerung-Paket und Compilieren
Kontrollausgaben 
Netzwerkkomponenten zusammen arbeiten lassen (Hard/-Software)
Wlan -> Teilproblem
Kamera -> Teilproblem Zeit


\section{Fazit}
1 Seite
\subsection{Zusammenfassung}

sehr knapper Zeitrahmen
hohe Anforderungen für diesen Zeitraum
interessantes Projekt
anspruchsvolles Projekt
viel über Linux und Ros und Hardware und Netzwerken gelernt
Spaß

\section{Ausblick}

Benutzung des Pi's als Hardware Nodes Lab
Benutzung des Pi's für kleine Roboterprojekte Lab
Immer mehr binary packages für ROS Pi verfügbar -> mehr Möglichkeiten
z.B.: GPIO Interface -> Knöpfe Kaffemaschine, Relays, alle möglichen Signale
auf Hardwareebene

\section{Eigene Leistung}

anstatt Linux PC wird Pi benutzt

\section{Quellen}
\begin{verbatim}

http://de.wikipedia.org/wiki/Raspberry_Pi
http://www.amazon.de/Raspberry-Pi-RBCA000-Mainboard-1176JZF-S/dp/B008PT4GGC
http://www.raspbian.org/
http://www.softwareok.de/?seite=faq-System-Allgemein&faq=13
http://de.wikipedia.org/wiki/ARM-Architektur
http://www.rn-wissen.de/index.php/Raspberry_PI:_GPIO

\end{verbatim}

\end{document}